% Dokumentum konfiguráció
\documentclass[fleqn,12pt]{article}
% --------------------------------------------------
% Dokumentum adatok
% --------------------------------------------------

% Intézmény és kar
\newcommand\INTEZMENY{Budapesti Műszaki és Gazdaságtudományi Egyetem}
\newcommand\KAR{Gépészmérnöki Kar}
\newcommand\LOGO{bme-logo.jpg} % logo fájl neve a figures mappában

% Dokumentum adatok
\newcommand\CIM{Grundlagen der Elektronik Hausaufgabe}
\newcommand\DATUM{2025/26/2}
\newcommand\NEV{Abos Gergő}
\newcommand\NEPTUN{C46M9M}
\newcommand\EMAIL{gergo.abos@gmail.com}
\newcommand\TANTARGY{Grundlagen der Elektronik}
\newcommand\TANTARGYKOD{BMEVIAUBSXA060-00}

\usepackage{tabularx}
% --------------------------------------------------
% Preambulum és egyéni parancsok betöltése
% --------------------------------------------------

\input{../_settings/00_preambulum.tex}
\input{../_settings/01_commands.tex}
\input{../_settings/02_mechTikz.tex}

% ==================================================
% Dokumentum kezdete
% ==================================================
\begin{document}
\textbf{sid1.}
\begin{align*}
&\text{A) 10V Gleichstrom} \\
&\text{B) 100V absolute Mittelwert Sinus} \\
&\text{C) 100V abs. Mittelwert, -U:50\%..+U:50\% Viereckspannung} \\
&\text{D) Dreieck Spannung, 25\% steigend, 75\% fallend, -100...100V}
\end{align*}
\begin{table}[htpb]
\centering
\renewcommand{\arraystretch}{1.3} % Etwas mehr Zeilenabstand für bessere Lesbarkeit
\begin{tabularx}{\textwidth}{|X|c|c|c|c|}
\hline
\textbf{Messwerk} & \textbf{Fall A} & \textbf{Fall B} & \textbf{Fall C} & \textbf{Fall D} \\ \hline
Durchschnittswert & 10V & 0V & & \\ \hline
Abs. Mittelwert & 10V & 100V & & \\ \hline
Effektivwert & 10V & 111,07V & & \\ \hline
Formfaktor & 1 & 1,1107 & & \\ \hline
Crestfaktor & 1 & 1,414 & & \\ \hline
Drehspul & 10V & 0V & & \\ \hline
Drehspul mit Gleichrichter & 11,1V & 111,07V & & \\ \hline
Dreheisen, Elektrodynamisch, Thermoumformer, Elektrostatisch & 10V & 111,07V & & \\ \hline
Digitale (einfache Typ mit einer Diode) in Gleichstrom Bereich & 10V & 50V & & \\ \hline
Digitale (einfache Typ mit einer Diode) in Wechselstrom Bereich & 0V & 111,07V & & \\ \hline
Digitale (einfache Typ mit einer Diode) in Wechselstrom Bereich, Anschlussleitungen vertauscht & 0V & 111,07V & & \\ \hline
Digitale (mit 2 Dioden, es erstellt den abs. Mittelwert) in Gleichstrom Bereich & 10V & 100V & & \\ \hline
Digitale (mit 2 Dioden, es erstellt den abs. Mittelwert) in Wechselstrom Bereich & 0V & 111,07V & & \\ \hline
Digitale (echt Effektivwert messende Typ) in Gleichstrom Bereich & 10V & 111,07V & & \\ \hline
Digitale (echt Effektivwert messende Typ) in Wechselstrom Bereich & 0V & 111,07V & & \\ \hline
\end{tabularx}
\end{table}
\pagebreak

\section{Fall A}
\begin{figure}[H]
\centering
\begin{tikzpicture}[>=Stealth]
\draw[->] (0,0) -- (10,0) node[above right]{$t$};
\draw[->] (0,0) -- (0,2) node[right]{$u(t)$};
\draw[thick, red] (0, 1.5) -- (9, 1.5) node[right]{$10V$};
\end{tikzpicture}
\end{figure}

Durchschnittswert, abs. Mittelwert, Effektivwert sind \textbf{10V}, da die Spannung konstant ist.

Formfaktor:
\begin{align*}
F = \frac{U_{eff}}{U_{abs}} = \frac{10}{10} = \mathbf{1}
\end{align*}

Crestfaktor:
\begin{align*}
C = \frac{U_{spitze}}{U_{eff}} = \frac{10}{10} = \mathbf{1}
\end{align*}

Drehspuhl zeigt \textbf{Durchschnittswert}, also es zeigt \textbf{10V}.

Drehspuhl mit Gleichrichter ist skaliert für ein Wechselstrom, also die Werte sind mit der Formfaktor einer Sinuswelle multipliziert. Es zeigt:
\begin{align*}
U \cdot F_{sin} = 10 \cdot 1,11 = \mathbf{11,1V}
\end{align*}
\end{document}
